\chapter{Corps et Somatisation : Quand le corps hurle}

\textit{La peau est une frontière, l'estomac est un alambic. Quand nous refusons de poser nos frontières par la parole, notre corps les dessine par la douleur. Ce qui ne s'exprime pas s'imprime.}

\section{Le Ventre, Mémoire du Silence}
Nous avons vu au Chapitre 2 l'impact du nerf vague sur nos viscères. Mais il faut aller plus loin. L'estomac est le premier théâtre de la soumission. Pour celui qui n'ose pas dire son désaccord, chaque repas pris dans la tension est une agression. Le système digestif, privé de l'énergie nécessaire car le cerveau est en mode "alerte sociale", se crispe. Les gastrites, les reflux et les colopathies fonctionnelles ne sont pas des maladies de l'organe, mais des maladies de la relation. Le corps tente de digérer l'indigeste : l'insulte que l'on a bue, le mépris que l'on a avalé.

\section{La Peau : La Frontière Blessée}
Pourquoi la peau réagit-elle si violemment au manque d'affirmation ? Parce qu'elle est notre interface avec le monde. Elle est ce qui nous sépare des autres. Lorsque nous ne savons pas dire "non", lorsque nous laissons les autres envahir notre territoire psychique, notre peau proteste. Le psoriasis, l'eczéma ou l'urticaire sont souvent des cris de la frontière. C'est une manière brutale pour le corps de dire : "Ne me touchez pas, je suis à bout".

\begin{NoteClinique}
\textbf{Le cas de Bruno, 50 ans.} \\
Bruno est d'une gentillesse extrême. Il accepte tout, subit tout avec un sourire navré. Mais depuis deux ans, Bruno souffre d'aphonies subites. Dès qu'une tension apparaît avec son épouse ou son directeur, sa voix s'éteint. Il n'a plus de son. Les examens ORL ne montrent rien. C'est son système nerveux qui, pour le protéger d'un conflit qu'il juge mortel, coupe littéralement le sifflet. Bruno ne parle plus car son corps a décidé que le silence absolu était sa seule chance de survie.
\end{NoteClinique}

\section{L'Immunité en Déroute}
Le stress de l'inhibition n'est pas un stress de combat, c'est un stress d'impuissance. Ce dernier est le plus destructeur pour le système immunitaire. Le haut taux de cortisol basal finit par inhiber la production de nos cellules défenseuses. En ne s'affirmant pas, on ne fragilise pas seulement ses relations, on fragilise son armure biologique. On devient perméable aux virus, comme on est devenu perméable aux volontés des autres.
